
%\documentclass[./main.tex(path of main file)]{subfiles}
\documentclass[10pt]{article}
\usepackage{amsmath}
\DeclareMathOperator*{\argmin}{arg\,min} % thin space, limits underneath in displays
\DeclareMathOperator*{\argmax}{arg\,max} % thin space, limits underneath in displays
\newtheorem{thm}{Theorem}
\usepackage{amssymb}
\usepackage{amsfonts}
\usepackage{mathrsfs}
\usepackage{bm}
\usepackage{indentfirst}
\setlength{\parindent}{0em}
\usepackage[margin=1in]{geometry}
\usepackage{graphicx}
\usepackage{setspace}
\doublespacing
\usepackage[flushleft]{threeparttable}
\usepackage{booktabs,caption}
\usepackage{float}
\usepackage[sort,comma]{natbib}
\usepackage[hidelinks]{hyperref}
\usepackage{booktabs}
\usepackage{multirow}

\usepackage{import}
\usepackage{xifthen}
\usepackage{pdfpages}
\usepackage{transparent}
\usepackage{subfiles}
\usepackage{blindtext}
\usepackage{glossaries}
\usepackage{pdflscape}
\usepackage{adjustbox}
\usepackage{rotating}
\usepackage{eurosym,geometry,ulem,subcaption,caption,color,sectsty,comment,footmisc,pdflscape,array}

% Format definition list
\newenvironment{mylist}[1]%
  {\begin{list}{}{\settowidth{\labelwidth}{\textbf{#1}}
   \setlength{\leftmargin}{\labelwidth}
   \addtolength{\leftmargin}{\labelsep}
   \renewcommand{\makelabel}[1]{\textbf{\hfill##1}}}}%
  {\end{list}}



\newcommand{\incfig}[1]{%
\def\svgwidth{\columnwidth}
\import{./figures/}{#1.pdf_tex}
}




\title{}
\author{}
\date{}


\begin{document}
\maketitle


\section{New Keynesian and Real Business Cycle model}

\subsection{What's being adopted by the NK from the RBC}
\begin{itemize}
\item 
The RBC model relies on dynamic, stochastic, general equilibrium framework.
The NK model has fully adopted these features.
\item 
The RBC approach takes the form of a set of equations that describe, in  a
highly aggregative manner: 1) the behavior of HHs, firms, and policymakers,
2) some market clearing and/or resource constraints, and 3) the evolution of
one or more exogenous variables that are the ultimate source of fluctuations 
in the economy.
{\textbf {More controversial}} may be the assumption, widely found in both 
RBC and NK models, that the behavior of HHs and firms is the outcome of an
optimization problem, solved under the assumption of rational expectations.
\end{itemize}


\subsection{What's being added, in the NK model, on top of RBC}
\begin{itemize}
\item 
		It introduces nominal variables explicitly: prices, wages, nominal 
		interest rates.
\item 
		It departs from the assumption of perfect competition in the goods
		market, allowing for positive price markups.
\item 
		It introduces nominal rigidities, generally using the formalism
		proposed by Calvo (1983), whereby only a constant fraction of 
		firms, drawn randomly from the population, are allowed to adjust
		the price of their good.
    The assumption of imperfect competition is often extended to the 
		labor market as well, with the introduction of wage rigidities (
		nominal or real).
\end{itemize}

\subsection{Key properties of the NK model}
\begin{itemize}
\item 
		Exogenous changes in the monetary policy (MP) have {\underline {nontrivial}} effects on
		real variables, not only on nominal ones.
\item 
		The economy's equilibrium response to any shock is not independent of the monetary policy rule
		in place, thus opening the door to a meaningful analysis of alternative monetary policy rules.
\end{itemize}




\bibliographystyle{plainnat}
\bibliography{my_bib}

\end{document}

